%HOMEWORK template for M 325K
\documentclass{article}
\headheight=7pt         \topmargin=14pt
\textheight=574pt       \textwidth=445pt
\oddsidemargin=18pt     \evensidemargin=18pt

%Begin package import

\usepackage{amsmath, amssymb, amsthm}
\usepackage{mathtools}
\usepackage{pdfpages}
\usepackage{tikz-cd}

%End package import
%Begin custom commands

\newcommand{\C}{\mathbb{C}}
\newcommand{\R}{\mathbb{R}}
\newcommand{\Q}{\mathbb{Q}}
\newcommand{\Z}{\mathbb{Z}}
\newcommand{\N}{\mathbb{N}}
\newcommand{\F}{\mathbb{F}}
\newcommand{\abs}[1]{\left\lvert#1\right\rvert}
\newcommand{\RM}[1]{\mathrm{#1}}
\newcommand{\BF}[1]{\textbf{#1}}
\newcommand{\e}{\epsilon}
\newcommand{\ceil}[1]{\lceil{#1}\rceil}
\newcommand{\floor}[1]{\lfloor{#1}\rfloor}
\newcommand{\rarr}[0]{\rightarrow}
\newcommand{\IM}[1]{\text{Im}(#1)}
\newcommand{\Hom}[0]{\text{Hom}}
\newcommand{\Pow}[1]{\text{Pow}(#1)}
\newcommand{\angles}[1]{\langle #1 \rangle}

%End custom commands
%Begin custom theorems

\newtheorem{innercustomgeneric}{\customgenericname}
\providecommand{\customgenericname}{}
\newcommand{\newcustomtheorem}[2]{%
\newenvironment{#1}[1]
{%
\renewcommand\customgenericname{#2}%
\renewcommand\theinnercustomgeneric{##1}%
\innercustomgeneric%
}
{\endinnercustomgeneric}
}

\newcustomtheorem{THRM}{Theorem}
\newcustomtheorem{LEM}{Lemma}
\newcustomtheorem{EX}{Exercise}
\newcustomtheorem{COR}{Corollary}
\newcustomtheorem{PROB}{Problem}
\newcustomtheorem{claim}{Claim}

\newenvironment{solution}
{\renewcommand\qedsymbol{$\blacksquare$}\begin{proof}[Solution]}
{\end{proof}}

\newenvironment{definition}
{\renewcommand\qedsymbol{$\blacksquare$}\begin{proof}[Definition]}
{\end{proof}}

%End custom theorems
\begin{document}

\title{Homework 7}
\author{Arham Lodha}%put your name here
\date{March 05, 2024}%enter the due date here
\maketitle

\begin{PROB}{1}\label{Problem 1.}
\end{PROB}
\begin{proof}
    Assume the contrary $x < 0$, then let $\epsilon = -x$. By definition of limit, $\exists K \in \N \ni \forall n \geq K, x_n - x < \epsilon \implies x_n < \epsilon + x \implies x_n < 0$, which contradicts the hypothesis. Thus by contradiction, $x > 0$.       

    $\forall \epsilon > 0$, 
    
    The proof is broken down into two cases $x = 0, x \neq 0$ 
    \begin{enumerate}
        \item \textbf{Case 1} $x = 0$: By the definition of a limit, $\exists K \in \N \ni \forall n \geq K, x_n - x < \epsilon^2$. Thus $x_n < \epsilon^2$. Since $x_n > 0$, thus $0 < x_n < \epsilon^2$. Thus $0 < \sqrt{x_n} < \sqrt{\epsilon^2} = \epsilon$, thus $-\epsilon < 0 < \sqrt{x_n} < \epsilon \implies \abs{\sqrt{x_n} - 0} < \epsilon \implies \lim {\sqrt{x_n}} = 0 = \sqrt{x}$
        \item \textbf{Case 2} $x \neq 0$: $\forall \epsilon > 0$. By the definition of a limit, $\exists K \in \N \ni \forall n \geq K, \abs{x_n - x} < x\epsilon$. $\forall n \geq K$.    
    
    
        \begin{align*}
            \abs{\sqrt{x_n}-\sqrt{x}} &< \abs{\sqrt{x_n}-\sqrt{x}} \frac{\abs{\sqrt{x_n} + \sqrt{x}}}{x} \text{ (since $\frac{\abs{\sqrt{x_n} + \sqrt{x}}}{x} > 0$) } \\
            &= \frac{\abs{x_n - x}}{x} \\
            &< \epsilon \text{ since $\abs{x_n - x} < x\epsilon$ } \\
        \end{align*} 
        
        Thus $\lim \sqrt{x_n} = \sqrt{x}$. 
    \end{enumerate}
    
    Thus $\left(\sqrt{x_n}\right)$ is convergent to $\sqrt{x}$  
\end{proof}

\bigskip

\begin{PROB}{2}\label{Problem 2.}
\end{PROB}
\begin{proof}
    % Since $L < 1$, $1 - L > 0$.   
    % By the definition of limit, $\exists K \in \N \ni \forall n \geq K, \frac{x_{n + 1}}{x_n} - L < 1 - L$. Thus, $\frac{x_{n + 1}}{x_n} < 1 \implies x_{n + 1} < x_n$. Thus the tail of $x$ is monotonically decreasing, or the sequence $y_n = x_{n + K - 1}$ is monotonically decreasing. Furthermore since $\forall n \in \N$, $0 < x_n$, by assumption. $x_n$ is bounded and thus $y_n$ is bounded. 
    
    % By the Monotonic Bounded Theorem, $y_n$ has a limit and by the uniqueness of the limit, $x_n$ has a limit and $\lim x_n = \lim y_n$.
    
    Limit of a nonnegative sequence is nonnegative, $x_n > 0 \implies L \geq 0$.  
    
    Let $r \in \R$ be a number such that $0 \leq L < r < 1$. Let $\epsilon = r - L$. 
    
    
    By definition of the limit, $\exists K \in \N \ni \forall n \geq K, \frac{x_{n + 1}}{x_n}  - L< \epsilon$. Thus $\frac{x_{n + 1}}{x_n} < \epsilon +L = r \implies x_{n + 1} < rx_n$. Moreover since $\forall m \in \N$, $0 < x_m$, $0 < x_{n + 1} < rx_n$. 
    
    Thus $0 < x_{n + 1} < rx_{n} < r^2 x_{n - 1} <\dots < r^{n - K + 1}x_{K}$.
    
    Let $C := x_K r^{-K}$.
    
    Thus $0 < rx_{n} < Cr^{n + 1} \implies 0 < x_n < Cr^n$.
    
    $\forall \epsilon_2 > 0$. Let $M := \frac{1}{r} - 1$, thus $r = \frac{1}{M + 1}$. By the archemedian property, $\exists N \in \N \ni N > \frac{1}{M \epsilon} \implies \epsilon > \frac{1}{NM}$. $\forall n \geq N$ 
    
    \begin{align*}
        |r^n  - 0| = |r^n| = r^n &= \frac{1}{(M + 1)^n} \\
        &= \frac{1}{\sum_{i = 0} {n \choose i} M^i} \text{ (Binomial Expansion)} \\
        &< \frac{1}{1 + Mn} \text{ (since $M > 0$ and $1 > 0$ all terms of the sum are positive thus $(1 + Mn) <(M + 1)^n $)}
        &< \frac{1}{Mn} \text{ (since $1 + Mn > Mn$)}
        &< \frac{1}{MN} \text{ (since $n \geq N$ )} \\
        &< \epsilon \text{ (since $\epsilon > \frac{1}{MN}$ )}
    \end{align*}

    Thus $\lim(r_n) = 0$. By Properties of Limit operation, since $C$ is a constant, $\lim(Cr_n)= C\lim(r_n) = 0$.
    
    Since $0 < x_n < Cr^n$, by squeeze theorem $\lim(x_n) = 0$.  
\end{proof}

\bigskip

\begin{PROB}{3}\label{Problem 3.}
\end{PROB}
\begin{proof}
    Base Case: $1 \leq x_1 = 1 < \sqrt{3} = x_2 < 2$
    
    Induction Hypothesis: Suppose for $n = k$, $1 \leq x_k < x_{k + 1} < 2$
    
    Inductive Step: Thus $3 \leq x_{k} + 2 < x_{k + 1} + 2 < 4$ and $1 < \sqrt{3} \leq \sqrt{x_{k} + 2} = x_{k + 1} < \sqrt{x_{k + 1} + 2} = x_{k + 2} < 2$. Thus $1 \leq x_k < x_{k + 1} < 2 \implies 1 \leq x_{k + 1} < x_{k + 2} < 2$. 
    
    Thus by induction, $\forall n \in \N$ $1 \leq x_n < x_{n + 1} < 2$. 
    
    Thus $x_n$ is bounded and monotonically increasing. By montonic convergence theorem, $x_n$ converges to a limit. 
    
    $X = x_n$ and $X_1 = x_{n + 1}$, since the $X$ converges $X_1$ must converge to the same value by the uniqueness of the limit, since it is the tail of $X$. 
    
    Thus, $\lim x_n = \lim x_{n + 1} = \lim \sqrt{x_n + 2}$, By problem 1, $\lim \sqrt{x_n + 2} = \sqrt{\lim{x_n + 2}}$, by properties of operations on limits, $\lim{x_n + 2} = \lim{x_n} + \lim{2}$ limit of a constant is a constant, thus $\lim{x_n + 2} = \lim{x_n} + 2$. Thus, $\lim{x_n} = \sqrt{\lim{x_n} + 2}$. Let $L = \lim{x_n}$, thus $L = \sqrt{L^2 + 2} \implies L^2 - L - 2 = 0$ and since $x_n \geq 1 > 0$, $L > 0$, thus $L = 2$. Thus $\lim x_n = 2$.           
\end{proof}

\bigskip

\begin{PROB}{4}\label{Problem 4.}
\end{PROB}
\begin{proof}
    $\forall n \in \N$, $y_{n + 1} - y_n = \sum_{k = 1}^{n + 1} \frac{1}{n - k + 1} - \sum_{k = 1}^{n} \frac{1}{n + k} = - \frac{1}{n + 1} + \frac{1}{2n + 1} + \frac{1}{2n + 2} = \frac{1}{2n + 1} - \frac{1}{2n + 2}$, since $2n + 1 < 2n + 2$, $\frac{1}{2n + 1} > \frac{1}{2n + 2} \implies y_{n + 1} - y_n = \frac{1}{2n + 1} - \frac{1}{2n + 2} > 0 \implies y_{n + 1} > y_n$. Thus $y_n$ is montonically increasing.
    
    $\forall n \in \N$, $\forall k \in [1, \dots, n]$, $n + k > n \implies \frac{1}{n + k} < \frac{1}{n}$. Thus  $y_n = \sum_{k = 1}^{n} \frac{1}{n + k}< \sum_{k = 1}^{n} \frac{1}{n} = \frac{n}{n} = 1$. Thus $y_n$ is bounded above by 1. 
    
    Since $y_n$ is montonically increasing and bounded above by 1, by montonic convergence theorem it is convergent.  
\end{proof}

\bigskip

\begin{PROB}{5a}\label{Problem 5a.}
\end{PROB}
\begin{proof}
    $a_n = (1 - (-1)^n + \frac{1}{n})$. 
    Let $\epsilon = 1$. $\forall N \in \N$, if $N$ is even let $n = N + 1$, if $N$ is odd, let $n = N$. Let $m = 2n$.           

    \begin{align*}
        \abs{a_{m} - a_n} &= \abs{1 - (-1)^{2n} + \frac{1}{2n} - 1 + (-1)^n - \frac{1}{n}} \\
        &= \abs{(-1)^n - (-1)^{2n} + \frac{1}{2n} - \frac{1}{n}} \\
        &= \abs{-1 - 1 - \frac{1}{2n}} \\
        &= \abs{-2 - \frac{1}{2n}} \\
        &= 2 + \frac{1}{2n} \\
        &> \epsilon \text{ since $2n > 0$ and thus $2 + \frac{1}{2n} > 0$ }
    \end{align*}

    Thus $a_n$ is not cauchy. By cauchy convergence theorem, and thus divergent (not convergent).  
\end{proof}

\begin{PROB}{5b}\label{Problem 5b.}
\end{PROB}
\begin{proof}
    $a_n = \sin(\frac{n\pi}{4})$

    Let $\epsilon = 1$. $\forall N \in \N$.
    
    Let $n = N + 6 - (N \mod 8)$ and let $m = N + 10 - (N \mod 8)$. Thus $n \mod 8 = 6$ and $m \mod 8 = 2$ thus $a_n = \sin(\frac{(N + 6 - N \mod 8) \pi}{4}) = \sin(6 \pi / 4) =  -1$ and $a_m = \sin(\frac{(N + 10 - N \mod 8) \pi}{4}) = \sin(10 \pi /4) = 1$.      
    
    \begin{align*}
        \abs{a_m - a_n} &= \abs{1 - (-1)} \\
        &= \abs{2} = 2 > \epsilon
    \end{align*}

    Thus $a_n$ is not cauchy. By cauchy convergence theorem, and thus divergent (not convergent).  
\end{proof}

\bigskip

\begin{PROB}{6a}\label{Problem 6a.}
\end{PROB}
\begin{proof}
    $\forall \epsilon > 0$. By the Archemedian Property, $\exists K \in \N \ni K > \frac{1}{\epsilon} \implies \epsilon > \frac{1}{K}$. $\forall n \geq K$,    
    
    \begin{align*}
        \abs{x_m - x_n} = \abs{\frac{m + 1}{m} - \frac{n + 1}{n}} &= \abs{\frac{n - m}{mn}} \\
        &= \frac{m - n}{nm} \\
        &= \frac{1}{n} - \frac{1}{m}  
    \end{align*}

    $\frac{1}{n} - \frac{1}{m} - \frac{1}{n} = -\frac{1}{m} < 0$, since $m > 0$. Thus $\frac{1}{n} - \frac{1}{m} < \frac{1}{n}$.
    
    Thus, 

    \begin{align*}
        \abs{x_m - x_n} &< \frac{1}{n} \\
        &< \frac{1}{K} \text{ (since $n \geq K$)} \\
        &< \epsilon 
    \end{align*}

    Thus $x_n = \frac{n + 1}{n}$ is Cauchy. 
\end{proof}

\begin{PROB}{6b}\label{Problem 6b.}
\end{PROB}
\begin{proof}
    \begin{align*}
        \abs{y_{m} - y_n} = \abs{\sum_{k = 1}^{m}\left(\frac{1}{2}\right)^k - \sum_{k = 1}^{n}\left(\frac{1}{2}\right)^k} &= \abs{\sum_{k = n + 1}^{m} \left(\frac{1}{2}\right)^k} \\
        &= \sum_{k = n + 1}^{m} \left(\frac{1}{2}\right)^k \\
        &< \sum_{k = n + 1}^{m} \left(\frac{1}{2}\right)^{n + 1} \\
        &= \frac{m - n}{2^{n + 1}}
    \end{align*}
\end{proof}

\begin{PROB}{6b}\label{Problem 6b.}
\end{PROB}
\begin{proof}
    Let $r = \frac{1}{2}$. 
    
    $\forall n \in \N$, $y_n = \sum_{i = 1}^{n} r^i$, $ry_n = r\sum_{i = 1}^{n}r^{i} = \sum_{i = 2}^{n + 1} r^{i}$. $y_n(1 - r) = y_n - ry_n = r - r^{n + 1}$. Thus $y_n = \frac{r - r^{n + 1}}{1 - r} = \frac{\frac{1}{2} - \frac{1}{2}^{n + 1}}{\frac{1}{2}} = 1 - \left(\frac{1}{2}\right)^n$. Thus $y_n = 1 - \left(\frac{1}{2}\right)^n$ 


    \textbf{Claim: $\forall n \in \N$, $r^n < 1$.}   

    Base Case: n = 1, $r^n = \frac{1}{2} < 1$

    Induction Hypothesis: Suppose $r^k < 1$

    Inductive Step: $r r^k < r^{k + 1} < r^k \implies r^{k + 1} < 1$. 
    
    Thus by induction, $\forall n \in \N$, $r^n < 1$.  
    
    $\forall N \in \N$, 
    \textbf{Claim: $\forall n \geq N$, $r^n < r^N$.}

    Base Case: n = 1, $r^{N + 1} - r^N = r^{N}(r - 1) < 0$ since $\frac{1}{2} < 1$.

    Induction Hypothesis: Suppose $r^{N + k} < r^N$

    Inductive Step: $r r^{N + k} = r^{k + 1 + N} < r^{N + 1} \implies r^{k + 1} < r^N$. 
    
    Thus by induction, $\forall n \geq N$, $r^n < r^N$. 
             

    $\forall \epsilon > 0$. For $\epsilon > 1$, let $K = 0$, thus $\left(\frac{1}{2}\right)^K = 1 < \epsilon$. For $\epsilon < 1$, let $K = \ceil{-\log_2(\epsilon)}$. Since $\ceil{-\log_2(\epsilon)} > -\log_2(\epsilon)$, $2^{\ceil{-\log_2(\epsilon)}} > 2^{-\log_2(\epsilon)} \implies 2^{-\ceil{-\log_2(\epsilon)}} < 2^{\log_2(\epsilon)}$  
    
    
    Thus $\left(\frac{1}{2}\right)^K = \left(\frac{1}{2}\right)^{\ceil{-\log_2(\epsilon)}} < \left(\frac{1}{2}\right)^{-log_2(\epsilon)} = \epsilon$. Thus $\left(\frac{1}{2}\right)^K < \epsilon$. $\forall m , n \geq K \ni m > n$.              
    
    \begin{align*}
        \abs{y_m - y_n} = \abs{\left(\frac{1}{2}\right)^n - \left(\frac{1}{2}\right)^m} &= \abs{\left(\frac{1}{2}\right)^n \left(1 - \left(\frac{1}{2}\right)^{n - m}\right)} \\
        &= \left(\frac{1}{2}\right)^n \left(1 - \left(\frac{1}{2}\right)^{n - m}\right) \\
        &< \left(\frac{1}{2}\right)^n \text{ (since $r^{n - m} < 1$)}\\ 
        &< \left(\frac{1}{2}\right)^K \text{ (since $n \geq N \implies r^n < r^K$ )} \\
        &< \epsilon
    \end{align*}
\end{proof}

\bigskip

\begin{PROB}{7}\label{Problem 7.}
\end{PROB}
\begin{proof}
    
    $\forall \epsilon > 0$. By the archemedian property, $\exists K  \in \N \ni K > \frac{1}{4\epsilon^2} \implies \epsilon > \frac{1}{2\sqrt{K}}$  

    \begin{align*}
        \abs{\abs{x_{n + 1} - x_n} - 0} &= \abs{\sqrt{n + 1} - \sqrt{n}} \\
        &= \abs{\sqrt{n + 1} - \sqrt{n}}\frac{\abs{\sqrt{n + 1} + \sqrt{n}}}{\abs{\sqrt{n + 1} + \sqrt{n}}} \\
        &= \frac{\abs{1}}{\abs{\sqrt{n + 1} + \sqrt{n}}} \\
        &= \frac{1}{\sqrt{n + 1} + \sqrt{n}} \text{ (since $\sqrt{n + 1} > 0$ and $\sqrt{n} > 0$)} \\
        &< \frac{1}{2\sqrt{n}} \text{ (since $\sqrt{n + 1} > \sqrt{n}$)} \\
        &< \frac{1}{2\sqrt{K}} \\
        &< \epsilon \text{ (since $\epsilon > \frac{1}{2\sqrt{K}}$)}
    \end{align*}

    Thus, $\lim |x_{n + 1} - x_n| = 0$. 

    Assume the contrary, $\exists n \in \N \ni \sqrt{n} < 1$, thus $\sqrt{n} \sqrt{n} = n < \sqrt{n} \implies n < 1$. This is a contradiction, because $\forall n \in \N$, $n \geq 1$. Thus by contradiction, $\sqrt{n} \geq 1$     

    Let $\epsilon = \frac{1}{2}$. $\forall M \in \N$, $n = M$ and $m = 3M$.
    
    
    
    \begin{align*}
        \abs{x_m - x_n} = \abs{\sqrt{m} - \sqrt{n}} &= \abs{\frac{m - n}{\sqrt{m} + \sqrt{n}}} \\
        &= \frac{m - n}{\sqrt{m} + \sqrt{n}} \text{ (since $m > n$ and $\sqrt{m} + \sqrt{n} > 0$)} \\
        &< \frac{m - n}{2\sqrt{n}} \text{ (since $m > n$ so $\sqrt{m} > \sqrt{n}$)} \\
        &= \frac{3M - M}{2\sqrt{M}} \\
        &= \sqrt{M} \\
        &\geq 1 > \frac{1}{2} = \epsilon
    \end{align*}

    Thus $x_n$ is not cauchy. 
\end{proof}


\end{document}